%% -------------------------------------------------------
%% Kapitel 8: Zusammenfassung
%% -------------------------------------------------------
\chapter{Zusammenfassung}
\label{chap:zusammenfassung}

\section{Erreichte Ziele}
\label{sec:erreichte-ziele}

Im Rahmen dieser Diplomarbeit wurden die beiden zentralen Komponenten von DoggoGO –
Backend und iOS-Applikation – konzipiert und implementiert.

Das \textbf{Backend} auf Basis von ASP.NET~Core~8 stellt eine vollständige RESTful~API
mit Registrierung, JWT-Authentifizierung, Hundeprofilverwaltung und Rassendatenverwaltung
bereit. Die Datenhaltung erfolgt über PostgreSQL mit EF~Core~9; das automatische
Breed-Seeding aus einer JSON-Datei macht die Rassen-Datenbank sofort nutzbar. Das
gesamte System ist mittels Docker~Compose containerisiert und damit plattformunabhängig
deploybar.

Die \textbf{iOS-App} bietet eine kartenbasierte Darstellung des Benutzerstandorts
(MapKit), sichere JWT-Authentifizierung mit Keychain-Persistenz, Hundeprofilverwaltung
sowie die Anzeige rassenspezifischer Rechtsvorschriften. Die MVVM-Architektur mit
SwiftUI ermöglicht eine klare Trennung von Darstellung und Logik.

Alle in Kapitel~\ref{chap:anforderungen} definierten funktionalen Anforderungen
(FA-B01 bis FA-B08 und FA-I01 bis FA-I08) wurden erfüllt. Die nicht-funktionalen
Anforderungen bezüglich Sicherheit (NF-01, NF-02), Datenschutz (NF-03) und
Portierbarkeit (NF-06) wurden ebenfalls umgesetzt.

\section{Suboptimale Entscheidungen}
\label{sec:suboptimal}

Im Rückblick lassen sich einige Entscheidungen identifizieren, die bei einer Neuentwicklung
anders getroffen würden:

\paragraph{Doppelter Foreign-Key in der Initialmigration}
Die automatisch generierte EF~Core-Migration \texttt{InitDatabase} enthielt einen
redundanten Foreign-Key \texttt{Legal\-Restrictions\-Id1} in der \texttt{Breeds}-Tabelle.
Dieser entstand durch eine nicht klar definierte \texttt{WithMany()}-Konfiguration und
musste manuell aus der Migration entfernt werden. Eine sorgfältigere Konfiguration in
\texttt{OnModelCreating} hätte diesen Fehler vermieden.

\paragraph{Fehlendes Con\-trol\-ler-Layer}
Das Backend wurde zunächst ohne dedizierte Con\-trol\-ler-Klas\-sen begonnen; Endpunkte wurden
direkt in \texttt{Program.cs} als Minimal~API definiert. Erst im Verlauf des Projekts
wurde auf den klassischen Con\-trol\-ler-An\-satz umgestellt, was zu einer Umstrukturierung
des Codes führte. Ein frühzeitiger Entschluss für einen der beiden Ansätze hätte Zeit
gespart.

\paragraph{Fehlende Input-Validierung}
In der aktuellen Version werden Eingaben (z.\,B.\ beim Anlegen eines Hundes) nicht
serverseitig mit \texttt{DataAnnotations} oder einem Validierungsframework validiert.
Ungültige Werte (leerer Name, negatives Alter) könnten so in die Datenbank gelangen.
Für eine produktive Anwendung ist eine umfassende Validierungsschicht zwingend erforderlich.

\paragraph{iOS-Simulator statt echtes Gerät}
Die iOS-App wurde ausschließlich im Xcode-Simulator getestet. GPS-Standortdaten wurden
simuliert. Tests auf einem physischen Gerät hätten möglicherweise Probleme mit der
Standortgenauigkeit oder Keychain-Berechtigungen aufgedeckt.

\section{Einschränkungen der Lösung}
\label{sec:einschraenkungen}

\begin{itemize}
  \item \textbf{Keine echten Geo-Daten:} Freilaufwiesen, Sackerlspender und
    Hundekot-Mülleimer sind im aktuellen Stand nicht in der Datenbank enthalten.
    Die Kartendarstellung zeigt ausschließlich den Benutzerstandort. Die Entitäten und
    API-Endpunkte für Standortdaten wurden geplant, konnten im Rahmen der Diplomarbeit
    aber nicht vollständig implementiert werden.
  \item \textbf{Nur iOS, kein Android:} Eine Android-App ist nicht Teil des Projekts.
    Die API-Architektur ist jedoch so gestaltet, dass ein Android-Client ohne
    Backend-Änderungen ergänzt werden könnte.
  \item \textbf{Keine Produktionsumgebung:} Das System wurde für Entwicklungs- und
    Demonstrationszwecke konfiguriert. Für einen produktiven Betrieb wären HTTPS-Terminierung,
    Secrets-Management (z.\,B.\ Vault) und ein Reverse-Proxy (Nginx) notwendig.
  \item \textbf{Regionale Begrenzung:} Die Rechtsvorschriften sind derzeit auf
    österreichisches Recht (Wiener Hundegesetz) beschränkt. Eine Ausweitung auf andere
    Bundesländer oder Länder erfordert weitere Datenerhebung und -pflege.
\end{itemize}

\section{Erweiterungsmöglichkeiten}
\label{sec:erweiterungen}

Das System bietet eine solide Basis für zahlreiche Erweiterungen:

\begin{description}
  \item[Standort-Daten (DOG3–DOG5, DOG17)]
    Import von Open-Data-Geodaten (z.\,B.\ der Stadt Wien) für Freilaufwiesen,
    Sackerlspender und Hundewaschstationen. Neue \texttt{Location}-Entität im Backend,
    REST-Endpunkte für Geo-Abfragen.
  \item[Geofencing-Benachrichtigungen (DOG11)]
    Push-Benachrichtigungen via APNs, wenn die Nutzerin / der Nutzer eine Zone mit
    besonderen Vorschriften betritt. Serverseitig über einen Background-Worker realisierbar.
  \item[Routenplanung (DOG14)]
    Integration von MapKit-Routing zur Navigation zur nächsten Freilaufwiese.
  \item[Favoriten (DOG15)]
    Persistenz von Lieblings-Standorten in der Datenbank.
  \item[Android-App]
    Da die REST-API plattformunabhängig ist, kann ein Android-Client (Kotlin / Jetpack
    Compose) ohne Backend-Änderungen ergänzt werden.
  \item[Community-Features]
    Nutzerbewertungen und Erfahrungsberichte zu Standorten (neue Entitäten \texttt{Review},
    \texttt{Report}).
\end{description}
